% ===================================================================
%  図表第 7 号 — 冬の村。— 春の農家。
%
%  Ceci n'est PAS une transcription : c'est une traduction, faite en
%  2026, en japonais d'aujourd'hui. Voir texto/ja/10-zuhyo-01.tex
%  pour la consigne, qui vaut pour les seize fichiers.
% ===================================================================

%%K t07-noto-1 noto
\VUnotes{143.2mm}{%
(*) 一度の食事の細目は図表第 6 号と第 13 号を見よ。%
}

%%K t07-apar-1 apar
\VUsaut{4.0mm}
\VUcentre{13.2pt}{90}{第二部}
\VUsaut{3.0mm}
\VUfilet{16mm}
\VUsaut{5.6mm}
\VUtitre{15.0pt}{60}{図表第 7 号}
\VUsaut{3.0mm}

%%K t07-tit-1 sub
\VUcentre{12.0pt}{0}{\textit{第一場。}}
\VUsaut{3.0mm}

%%K t07-tit-2 sub
\VUpk{10.2pt}{40}{冬の村。}
\VUsaut{4.0mm}

%%K t07-c1-01-1 p
1. --- 新年の休みに、私は父と新庄へ行った。小さな村で、父はそこに用が
あったのである。

%%K t07-c1-02-1 p
2. --- 我々が乗ったのは町とその村とのあいだを通う\VUgras{駅馬車}
\textsuperscript{(11)} であった。しかしたいそう寒く、これほど乗り心地のわるい車
での道中は少しも快くなかった。

%%K t07-c1-03-1 p
3. --- だから\VUgras{御者}が広場で、\textit{白熊}の\VUgras{宿屋}
\textsuperscript{(2)} の前に車を止めたのを見たときは、まことにうれしかった。
\VUgras{宿の主} \textsuperscript{(4)} が戸口で我々を待っていて、静かに自分の
\VUgras{煙管}をふかしていた。

%%K t07-c1-03-2 p
主は中へ入れと言い、炉でぱちぱち鳴る葡萄の枝の火の前に坐るのに、私は二度
言われるまでもなかった。そのときはもう、あの身を切る\VUgras{北風}も気に
ならなかった。その風は古い\VUgras{看板} \textsuperscript{(3)} をきしませ、煙突から
出る\VUgras{煙} \textsuperscript{(1)} を黒い渦にして運び去っていたのである。

%%K t07-c1-04-1 p
4. --- ちょうど食事の時刻であった。もう待たずに、出された質素だがうまい
食事をしっかりといただいた (*)。

%%K t07-tit-3 sub
\VUpk{10.2pt}{40}{幾つかの訪問。}
\VUsaut{3.0mm}

%%K t07-c1-05-1 p
5. --- 食後、私は保険の代理をしている父について、その顧客を訪ねに行った。
まず宿屋の隣に住む\VUgras{鍛冶} \textsuperscript{(5)} のところへ行った。馬に蹄鉄を
打っているところであった。私はその手伝の力に感心した。馬の\VUgras{蹄}を自分の
革の前掛にしっかり当てていて、そのあいだ鍛冶が大槌で\VUgras{蹄鉄}
\textsuperscript{(6)} を打ちつけていたのである。

%%K t07-c1-06-1 p
6. --- 次に村の\VUgras{靴屋} \textsuperscript{(7)} 老ジャックのところへ行った。
看板には立派な\VUgras{長靴}が描いてあるのに、履きつぶした古靴の底を張り替えて
いるだけであった。

%%K t07-c1-06-2 p
老人は笑って我々に言った。「町では私は『靴作り』でしたが、客が一人もあり
ませんでした。ここでは『靴直し』で、仕事が余るほどあります。」

%%K t07-c1-07-1 p
7. --- 老人は隣の\VUgras{髪結} \textsuperscript{(8)} のことを話した。その男の客は
誰も満足していない。彼の\VUgras{剃刀}はいつも刃こぼれし、彼の\VUgras{鋏}は
一度もよく切れたためしがない。しかし村にはその男しか髪結がいないので、人々は
その手で髭を剃らせ髪を刈らせるほかない。おまけに吝嗇で、口をきくのも面倒な男
である。

%%K t07-c1-08-1 p
8. --- 幸いほかの\VUgras{隣人}はその男のようではない。たとえば\VUgras{車大工}
\textsuperscript{(9)} はよい人で、働き者で世話好きである。車の直しを頼むのは気持
のよいことである。その仕事はいつもきちんと片づいている。

%%K t07-c1-09-1 p
9. --- 老ジャックは\VUgras{馬具屋} \textsuperscript{(10)} にも文句がなかった。
仕事の立てこむときにはときどき\VUgras{馬具}を縫うのを手伝ってやるのである。

%%K t07-c1-09-2 p
父が老人に家族のことを尋ねた。「みな達者です。孫娘は\VUgras{学校}
\textsuperscript{(12)} におります。\VUgras{女の先生} \textsuperscript{(13)} はあの子を
たいそう気に入っておいでで、よい\VUgras{生徒} \textsuperscript{(14)} です。うちの
いたずら息子とは違います。あれは遊ぶことしか頭にありません。\VUgras{先生}
\textsuperscript{(15)} もどうにもならず、何一つ教えこめません。」

%%K t07-tit-4 sub
\VUsaut{3.0mm}
\VUpk{10.2pt}{40}{村の噂。}
\VUsaut{3.0mm}

%%K t07-c1-10-1 p
10. --- それから老人は村の便りを聞かせてくれた。新しい学校が建つのだという。
\VUgras{村役場}のはとうに狭くなっていたからである。これはたやすいことで、
\VUgras{粉屋}の庭を一部買えばよいだけである。それはあの気の毒な男にはかえって
よい話であった。寒くなってからというもの、\VUgras{水車小屋} \textsuperscript{(16)}
の\VUgras{石臼}は麦を挽けずにいる。\VUgras{水車} \textsuperscript{(17)} が
\VUgras{氷} \textsuperscript{(25)} に閉じこめられているからで、その氷は
\VUgras{小川} \textsuperscript{(23)} のものである。

%%K t07-c1-11-1 p
11. --- 「建物といえば、新しい\VUgras{教会} \textsuperscript{(18)} を御覧に
なりましたか。雪をかぶってたいそう美しゅうございます。\VUgras{烏}
\textsuperscript{(19)} は自分の古い鐘楼が分からなくなって、しょんぼり
\VUgras{墓地} \textsuperscript{(20)} の大きな木に止まっております。」

%%K t07-c1-12-1 p
12. --- 墓地の話が出ると、靴屋は父の知っていた人で、この一年に亡くなった
人々を思い出した。「旦那さま、\VUgras{墓} \textsuperscript{(21)} がまたどれほど
増えたことでしょう、あの\VUgras{糸杉} \textsuperscript{(22)} の下に。うちの老いた
公証人も亡くなりました。村じゅうがその\VUgras{葬い}に出ました。」

%%K t07-c1-13-1 p
13. --- 「あそこにその跡取りが小川で氷すべりをしております。さきほども
おいでになって、切れた\VUgras{氷すべり靴} \textsuperscript{(26)} の紐を直せと
おっしゃいました。」

%%K t07-c1-14-1 p
14. --- 「それにあそこ、小川の岸にいるのが新しい\VUgras{村の番人}
\textsuperscript{(27)} でございます。もとは憲兵で、村のいたずら小僧の天敵です。
その\VUgras{脚絆} \textsuperscript{(29)} と\VUgras{制帽} \textsuperscript{(28)} を
見るなり逃げてしまいます。」

%%K t07-c1-15-1 p
15. --- 「おや、老ジョゼフがパン屋に\VUgras{薪束の車} \textsuperscript{(30)} を
運んでまいります。--- なるほど、と父が言った、あの\VUgras{騾馬}
\textsuperscript{(31)} の一組は見覚えがある。あの男に話があるのだ、ちょうどよい。
ではまた、老ジャック。」

%%K t07-c1-16-1 p
16. --- ちょうど出ようとしたとき、村の子供たちが学校から放たれて
\VUgras{氷の滑り道} \textsuperscript{(33)} へ押し寄せた。\VUgras{転んで}
\textsuperscript{(34)} 手足を折る危険もかまわずにである。その一人が車大工の
ところから自分の\VUgras{橇} \textsuperscript{(32)} を取ってきて、仲間の一人を
乗せ、駆け出した。

%%K t07-c1-17-1 p
17. --- ほかの子供たちは\VUgras{雪の山} \textsuperscript{(37)} へ突進した。
\VUgras{橋} \textsuperscript{(40)} のそばである。そして前の日に作った
\VUgras{雪だるま} \textsuperscript{(39)} を撃ち始めた。その雪だるまには帽子を
かぶせ、煙管をくわえさせ、肩に学校鞄まで掛けてやってあった。

%%K t07-c1-18-1 p
18. --- 氷のような風も寒さも子供たちは気にかけない。多くは\VUgras{襟巻}
\textsuperscript{(35)} さえ持っていない。\VUgras{雪合戦} \textsuperscript{(38)} の
あとで暖まるには、熱い\VUgras{栗} \textsuperscript{(42)} を一握り買えばそれで足りる。
老ジャクリーヌという\VUgras{物売}から買うのである。その老女は薄すぎる
\VUgras{肩掛} \textsuperscript{(43)} にくるまって寒さに震えている。

%%K t07-tit-5 sub
\VUsaut{3.0mm}
\VUcentre{12.0pt}{0}{\textit{第二場。}}
\VUsaut{3.0mm}

%%K t07-tit-6 sub
\VUpk{10.2pt}{40}{春の農家。}
\VUsaut{4.0mm}

%%K t07-c2-01-1 p
1. --- 冬にもそれなりの楽しみはあるが、それでも我々は\VUgras{春}の帰りを
心から喜んで迎える。

%%K t07-c2-01-2 p
五月の夕暮の農家の庭は、なんとうれしい眺めであろう。

%%K t07-c2-02-1 p
2. --- 地平では\VUgras{太陽} \textsuperscript{(1)} がゆっくり沈み、\VUgras{野}
\textsuperscript{(3)} をその紅の\VUgras{光} \textsuperscript{(2)} に浸している。働き者
の\VUgras{耕す人} \textsuperscript{(6)} は自分の\VUgras{畑} \textsuperscript{(4)} を
急いで耕している。

%%K t07-c2-02-2 p
自分の\VUgras{犂} \textsuperscript{(5)} で --- それは丈夫な\VUgras{馬}
\textsuperscript{(7)} が引いている --- \VUgras{畝} \textsuperscript{(8)} を切って
ゆく。そこへ\VUgras{種蒔く人} \textsuperscript{(9)} が一握りずつ\VUgras{種}を
まくのである。やがて黄金の穂となる種である。

%%K t07-c2-03-1 p
3. --- \VUgras{草刈} \textsuperscript{(11)} たちは鎌で草地の草を刈り、干す者は
\VUgras{干草} \textsuperscript{(10)} を\VUgras{草叉} \textsuperscript{(12)} と熊手で
返して乾かし、いまは帰るところである。

%%K t07-c2-03-2 p
ある者は牛の引く重い車の前を歩き、肩に道具を担いでいる。ほかの者は、もっと
横着で、香りのよい干草の上にゆったりと寝そべっている。

%%K t07-c2-04-1 p
4. --- 通りすがりに彼らは\VUgras{庭師} \textsuperscript{(16)} に親しげに挨拶する。
庭師は\VUgras{果樹園} \textsuperscript{(13)} で\VUgras{果樹}を刈りこみ、
\VUgras{梨の木} \textsuperscript{(14)} に接木をしている。\VUgras{李の木}
\textsuperscript{(15)} は花ざかりで、よく肥をやった\VUgras{菜園} \textsuperscript{(17)}
では野菜も甘藍も萵苣もすでによく育っている。

%%K t07-c2-04-2 p
低い塀の上では美しい\VUgras{孔雀} \textsuperscript{(18)} が尾を広げ、見事な羽を
見せびらかしている。

%%K t07-tit-7 sub
\VUpk{10.2pt}{40}{農家の庭。}
\VUsaut{3.0mm}

%%K t07-c2-05-1 p
5. --- \VUgras{馬屋} \textsuperscript{(19)} の戸は開いていて、飼葉桶と
\VUgras{敷藁} \textsuperscript{(23)} が見える。それは\VUgras{牝馬}
\textsuperscript{(21)} のもので、\VUgras{馬丁} \textsuperscript{(20)} がいま
車から外したばかりである。\VUgras{仔馬} \textsuperscript{(22)} は長い脚がなんとも
おかしく、跳ねながら母のあとを追う。

%%K t07-c2-06-1 p
6. --- \VUgras{井戸} \textsuperscript{(29)} のそばへ\VUgras{牝牛}
\textsuperscript{(25)} が水を飲みに来る。飲むのは水槽がわりの木の\VUgras{桶}
\textsuperscript{(26)} である。\VUgras{仔牛} \textsuperscript{(24)} はその隙に乳を
飲み、\VUgras{牡牛} \textsuperscript{(27)} は飼葉のよい匂いを嗅いで、うれしげに
鳴き、力強い\VUgras{角} \textsuperscript{(28)} のついた頭を誇らしげに上げる。

%%K t07-c2-07-1 p
7. --- 誰もが働いている。\VUgras{下女} \textsuperscript{(32)} が井戸の\VUgras{綱}
\textsuperscript{(31)} を引いている。\VUgras{手桶} \textsuperscript{(30)} で水を汲み、
家畜に飲ませるのである。\VUgras{穀倉} \textsuperscript{(33)} のそばでは一人の男が
散った藁を掻き集め、\VUgras{農家} \textsuperscript{(34)} の戸の前では老いた
\VUgras{糸紡ぎ女} \textsuperscript{(35)} が糸車と自分の\VUgras{紡錘}を前に、
昔ながらに\VUgras{亜麻}や\VUgras{大麻}を紡いでいる。

%%K t07-tit-8 sub
\VUsaut{3.0mm}
\VUpk{10.2pt}{40}{農家の主。}
\VUsaut{3.0mm}

%%K t07-c2-08-1 p
8. --- \VUgras{農家の主} \textsuperscript{(36)} がこの仕事すべてを差配し、下男に
指図する。\VUgras{馬鍬} \textsuperscript{(40)} と\VUgras{十字鍬} \textsuperscript{(39)}
を小屋にしまえと言いつけている。その二つは\VUgras{小屋} \textsuperscript{(38)} の
そばに置きっぱなしで、その小屋は\VUgras{犬} \textsuperscript{(37)} のものである。

%%K t07-c2-09-1 p
9. --- その声に驚いて\VUgras{鳩} \textsuperscript{(41)} が一羽、力いっぱい
\VUgras{鳩小屋} \textsuperscript{(42)} へ飛んでゆき、\VUgras{雌鶏}
\textsuperscript{(43)} たちも驚いて急いで\VUgras{鶏小屋} \textsuperscript{(44)} へ
戻ってゆく。

%%K t07-c2-10-1 p
10. --- \VUgras{羊小屋} \textsuperscript{(48)} の前では、老いた\VUgras{羊飼}
\textsuperscript{(45)} が自分の\VUgras{群} \textsuperscript{(49)} を連れ戻している。
自分の\VUgras{杖} \textsuperscript{(46)} を持っていて、その杖は色のあせた自分の
\VUgras{上っ張} \textsuperscript{(47)} と同じく決して手を離れない。彼は囲いへ
\VUgras{牡羊} \textsuperscript{(50)}、\VUgras{牝羊} \textsuperscript{(53)}、
\VUgras{仔羊} \textsuperscript{(54)}、\VUgras{山羊} \textsuperscript{(51)}、
\VUgras{仔山羊} \textsuperscript{(52)} を追いこむ。忠実な\VUgras{犬}
\textsuperscript{(55)} アルゴスが一番後ろを行き、吠えて遅れる者を追い立てる。

%%K t07-c2-11-1 p
11. --- この騒がしさもあの善良な\VUgras{乳牛} \textsuperscript{(56)} の落ち着きを
乱さない。牛は自分の\VUgras{鈴} \textsuperscript{(57)} を鳴らしながら、\VUgras{牛舎}
\textsuperscript{(58)} の戸の前で乳を搾らせている。

%%K t07-c2-12-1 p
12. --- 牛は自分が乳を出さねばならぬと心得ているようである。\VUgras{農家の妻}
\textsuperscript{(60)} が\VUgras{攪乳器} \textsuperscript{(61)} で\VUgras{バター}を
作れるように、またあの上等の\VUgras{チーズ} \textsuperscript{(64)} を作れるように
である。チーズは\VUgras{桶} \textsuperscript{(63)} に渡した板の上で水を切って
いる。

%%K t07-c2-13-1 p
13. --- \VUgras{乳室} \textsuperscript{(59)} じゅうが\VUgras{農家の下男}
\textsuperscript{(65)} のおかげでたいそうきれいに保たれている。彼はいま手に提げた
大きな桶で\VUgras{牛乳の缶} \textsuperscript{(62)} を洗ったところである。

%%K t07-tit-9 sub
\VUsaut{3.0mm}
\VUpk{10.2pt}{40}{鶏の庭。}
\VUsaut{3.0mm}

%%K t07-c2-14-1 p
14. --- 重い木靴をはいてよろよろ歩くので、\VUgras{七面鳥} \textsuperscript{(68)}、
\VUgras{牝鴨} \textsuperscript{(66)}、\VUgras{牡鴨} \textsuperscript{(67)}、
\VUgras{鵞鳥} \textsuperscript{(69)} が驚いて逃げる。みな\VUgras{嘴}
\textsuperscript{(70)} を大きく開けて不安げに鳴きたてる。

%%K t07-c2-14-2 p
\VUgras{雄鶏} \textsuperscript{(71)} はもっと気が荒く、赤い\VUgras{鶏冠}
\textsuperscript{(72)} を立て、自分の\VUgras{尾} \textsuperscript{(73)} の羽を振るい、
高らかに「こけこっこう」と鳴く。

%%K t07-c2-15-1 p
15. --- 一羽の\VUgras{母鶏} \textsuperscript{(74)} が\VUgras{糞の山}
\textsuperscript{(81)} をつついている。自分の\VUgras{雛} \textsuperscript{(75)} を
連れてである。\VUgras{仔豚} \textsuperscript{(78)} が母のほうへ駆けてゆく。豚小屋の
そばに見える大きな\VUgras{牝豚} \textsuperscript{(77)} である。

%%K t07-c2-16-1 p
16. --- 籠の中では\VUgras{兎} \textsuperscript{(79)} が、農家の主の娘が持ってきた
やわらかい\VUgras{草} \textsuperscript{(80)} を勢いよく食べている。ジャネットは
自分の兎をたいそう可愛がっていて、毎日、鶏小屋から新しい\VUgras{卵}
\textsuperscript{(76)} を取ってきたあと、この小さな友だちに会いに来る。
\VUsaut{6.0mm}
\VUfilet{22mm}
